\section{Implement Trie (Insert, Search, StartsWith)}
\label{sec:trie-implement}

\problemstatement{Implement a trie supporting $\textit{insert}(word)$,
$\textit{search}(word)$ (returns true only if $word$ was previously
inserted as a \emph{complete} word), and $\textit{startsWith}(prefix)$
(returns true if any inserted word begins with $prefix$).}

\begin{patternbox}
This is the foundational trie, with no special twist -- the rest of this
chapter builds directly on the node structure and walk logic derived
here. The one subtlety worth flagging immediately: $\textit{search}$ and
$\textit{startsWith}$ use the \emph{same} walk, differing only in what
they check once the walk finishes.
\end{patternbox}

\subsection*{Understanding the problem carefully}

Each trie node represents one character position along some prefix, and
holds up to $26$ child pointers (one per lowercase letter, say), plus a
boolean flag marking whether a complete inserted word \emph{ends} exactly
at this node. Walking a string through the trie, one character at a
time, either follows existing child pointers (if a previous word shares
this prefix) or creates new nodes (during insertion, for the first word
to diverge here).

\begin{ideabox}[One Walk, Three Different Endings]
$\textit{insert}(word)$: starting at the root, for each character,
create the corresponding child if it does not yet exist, then move into
it; after the last character, mark the final node's end-of-word flag.
$\textit{search}(word)$: walk the same way, but if any required child is
missing, return false immediately; after the last character, return
whether the final node's end-of-word flag is set (not just whether the
path exists -- $word$ might only be a \emph{prefix} of something
inserted, not itself a complete inserted word). $\textit{startsWith}
(prefix)$: identical walk, but return true as soon as the path exists in
full, regardless of the end-of-word flag.
\end{ideabox}

\subsection*{Why search must check the end-of-word flag but startsWith must not}

This is the one place a careless implementation typically goes wrong: if
only ``apple'' has been inserted, then $\textit{search}(\text{"app"})$
must return \texttt{false} (``app'' was never inserted as a complete
word, even though it is a valid path in the trie), while
$\textit{startsWith}(\text{"app"})$ must return \texttt{true} (something
inserted does begin with ``app''). The trie's path existing is necessary
but not sufficient for $\textit{search}$ -- it additionally needs the
explicit end-of-word marker, which is exactly the piece of information
$\textit{startsWith}$ does not need to check at all.

\subsection*{Algorithm}

\begin{enumerate}
  \item \textsc{Insert}($word$): $\textit{node} \gets \textit{root}$.
        For each character $c$ in $word$: if $\textit{node}$ has no
        child $c$, create one; move $\textit{node}$ into child $c$.
        Mark $\textit{node}.\textit{isEnd} \gets \text{true}$.
  \item \textsc{Search}($word$): walk as above; if any required child
        is missing, return false; otherwise return
        $\textit{node}.\textit{isEnd}$.
  \item \textsc{StartsWith}($prefix$): walk as above; return true if the
        full walk succeeds (regardless of $\textit{isEnd}$), false if any
        child is missing.
\end{enumerate}

\subsection*{Dry run}

\dryrun{Insert \texttt{"apple"}. Query \textit{search}(\texttt{"app"}),
\textit{startsWith}(\texttt{"app"}), \textit{search}(\texttt{"apple"}).}

\begin{itemize}
  \item Insert \texttt{"apple"}: creates a path of $5$ nodes,
        $a\to p\to p\to l\to e$, marking only the final ($e$) node's
        $\textit{isEnd}$ flag.
  \item $\textit{search}(\texttt{"app"})$: path $a\to p\to p$ exists,
        but that node's $\textit{isEnd}$ is false (only the $e$ node at
        the end of ``apple'' has it set). Returns \texttt{false}.
  \item $\textit{startsWith}(\texttt{"app"})$: path exists, so returns
        \texttt{true}, regardless of the flag.
  \item $\textit{search}(\texttt{"apple"})$: path exists and the final
        node's $\textit{isEnd}$ is true. Returns \texttt{true}.
\end{itemize}

\subsection*{C++ implementation}

\begin{lstlisting}
#include <bits/stdc++.h>
using namespace std;

struct TrieNode {
    TrieNode* children[26] = {nullptr};
    bool isEnd = false;
};

class Trie {
    TrieNode* root;

public:
    Trie() { root = new TrieNode(); }

    void insert(string word) {
        TrieNode* node = root;
        for (char c : word) {
            int idx = c - 'a';
            if (!node->children[idx]) node->children[idx] = new TrieNode();
            node = node->children[idx];
        }
        node->isEnd = true;
    }

    bool search(string word) {
        TrieNode* node = root;
        for (char c : word) {
            int idx = c - 'a';
            if (!node->children[idx]) return false;
            node = node->children[idx];
        }
        return node->isEnd;
    }

    bool startsWith(string prefix) {
        TrieNode* node = root;
        for (char c : prefix) {
            int idx = c - 'a';
            if (!node->children[idx]) return false;
            node = node->children[idx];
        }
        return true;
    }
};

int main() {
    Trie trie;
    trie.insert("apple");
    cout << trie.search("app") << endl;       // 0 (false)
    cout << trie.startsWith("app") << endl;   // 1 (true)
    cout << trie.search("apple") << endl;     // 1 (true)
    return 0;
}
\end{lstlisting}

\begin{complexitybox}
\textbf{Time Complexity.} $O(L)$ for each operation, where $L$ is the
length of the word/prefix involved -- independent of how many other
words are stored in the trie.
\textbf{Space Complexity.} $O(\Sigma \cdot N)$ in the worst case, where
$N$ is the total number of characters inserted across all words and
$\Sigma$ is the alphabet size (each node reserves space for all $26$
possible children, whether used or not); shared prefixes reduce this in
practice.
\end{complexitybox}

\begin{takeawaybox}
The entire chapter's machinery is this one node structure (a fixed-size
child array plus an end marker) and this one walk (follow or create
children, one character at a time). Every later section either adds more
annotations to each node (counts, in
Section~\ref{sec:trie-implement-2}) or changes what ``character'' means
(bits, in Sections~\ref{sec:trie-max-xor}--\ref{sec:trie-max-xor-query}).
\end{takeawaybox}
